\documentclass[11pt,a4paper]{article}
\input{preamble}

\pagestyle{fancy}
\fancyhf{}
\fancyhead[L]{\small\color{lightbrown}\emoji{graduation-cap} RL 틱택토 활용가이드}
\fancyhead[R]{\small\color{lightbrown}v1.0.0.0}
\fancyfoot[C]{\thepage}
\renewcommand{\headrulewidth}{0.4pt}
\setlength{\headheight}{14pt}

\begin{document}

\begin{titlepage}
\centering
\vspace*{3cm}
{\fontsize{48pt}{50pt}\selectfont \emoji{graduation-cap}}
\vspace{1cm}

{\Huge\bfseries\color{mainblue} RL 틱택토}
\vspace{0.3cm}

{\LARGE\color{darkbrown} 활용가이드}
\vspace{1.5cm}

{\large 교육용 Q-Learning 시뮬레이터의 수업 적용 가이드}
\vspace{3cm}

\begin{tabular}{rl}
\textbf{문서 버전} & v1.0.0.0 \\
\textbf{시뮬레이터} & rl-tictactoe-v3.html (v3) \\
\textbf{작성일} & 2026-02-14 \\
\end{tabular}
\vfill
{\small\color{lightbrown} 청주교육대학교 — 컴퓨팅사고와 AI교육}
\end{titlepage}

\tableofcontents
\newpage

% ============================================================
\section{이 문서의 목적}
% ============================================================

\subsection{활용가이드란?}

이 문서는 \textbf{RL 틱택토 시뮬레이터를 초등교육 현장에서 활용하는 방법}을 안내합니다. 예비 초등교사가 ``컴퓨팅 사고와 AI 교육'' 수업에서 강화학습 개념을 가르칠 때 바로 사용할 수 있는 수업 시나리오, 활동지, 질문 모음을 제공합니다.

\subsection{관련 문서 읽기 순서}

\begin{center}
\renewcommand{\arraystretch}{1.4}
\begin{tabular}{|c|l|l|c|}
\hline
\rowcolor{tableheader}
\textbf{순서} & \textbf{문서} & \textbf{목적} & \textbf{소요 시간} \\
\hline
1 & \textbf{이론해설서} & Q-Learning 원리 이해 & 20분 \\
\hline
2 & \textbf{사용설명서} & 시뮬레이터 전체 기능 파악 & 15분 \\
\hline
3 & \textbf{활용가이드 (본 문서)} & 수업 적용 방법 & 10분 \\
\hline
4 & 퀵가이드 & 수업 중 빠른 참조 & 필요 시 \\
\hline
\end{tabular}
\end{center}

\subsection{수업 전 준비 사항}

\begin{center}
\renewcommand{\arraystretch}{1.4}
\begin{tabular}{|l|p{9cm}|}
\hline
\rowcolor{tableheader}
\textbf{항목} & \textbf{준비 내용} \\
\hline
파일 & \texttt{rl-tictactoe-v3.html}을 학생 PC에 배포 \\
\hline
브라우저 & Chrome, Edge, 또는 Firefox (최신 버전) \\
\hline
사전 학습 & 교사가 시뮬레이터를 10분 이상 직접 사용해 볼 것 \\
\hline
인쇄물 & 활동지 3종 인쇄 (§3 참조) \\
\hline
선택 & 이론해설서 §1\textasciitilde2 요약 슬라이드 \\
\hline
\end{tabular}
\end{center}

% ============================================================
\section{수업 시나리오 (45분)}
% ============================================================

\subsection{수업 개요}

\begin{center}
\renewcommand{\arraystretch}{1.4}
\begin{tabular}{|l|p{9cm}|}
\hline
\rowcolor{tableheader}
\textbf{항목} & \textbf{내용} \\
\hline
대상 & 초등 예비교사 (대학 2\textasciitilde4학년) \\
\hline
과목 & 컴퓨팅 사고와 AI 교육 \\
\hline
주제 & 강화학습의 원리 — AI가 게임을 배우는 방법 \\
\hline
학습 목표 & \emoji{check-mark-button} 강화학습의 3요소를 설명할 수 있다 \newline \emoji{check-mark-button} 파라미터 변경이 AI 학습에 미치는 영향을 실험으로 확인할 수 있다 \newline \emoji{check-mark-button} AI의 학습 과정을 관찰하고 교육적 시사점을 도출할 수 있다 \\
\hline
준비물 & PC (1인 1대), rl-tictactoe-v3.html, 활동지 \\
\hline
\end{tabular}
\end{center}

\subsection{수업 흐름 (5단계)}

\begin{center}
\renewcommand{\arraystretch}{1.4}
\begin{tabular}{|c|c|p{3cm}|p{3cm}|p{2.5cm}|}
\hline
\rowcolor{tableheader}
\textbf{단계} & \textbf{시간} & \textbf{활동} & \textbf{시뮬레이터 기능} & \textbf{교사 역할} \\
\hline
\emoji{clapper-board} 도입 & 0\textasciitilde5분 & 동기 유발 & — & 시연 \\
\hline
\emoji{books} 개념 & 5\textasciitilde13분 & 핵심 원리 설명 & 도움말 모달 & 설명 + 비유 \\
\hline
\emoji{magnifying-glass-tilted-left} 탐색 & 13\textasciitilde23분 & 교사 주도 관찰 & 대시보드 학습 & 안내 + 질문 \\
\hline
\emoji{test-tube} 실습 & 23\textasciitilde37분 & 학생 주도 실험 & 파라미터, 대전 & 순회 지도 \\
\hline
\emoji{memo} 정리 & 37\textasciitilde45분 & 성찰 및 공유 & 학습 보고서 & 정리 + 피드백 \\
\hline
\end{tabular}
\end{center}

\subsection{단계별 상세}

\subsubsection{\emoji{clapper-board} 도입 (5분) — ``AI도 게임을 배울 수 있을까?''}

교사가 ``짝꿍과 틱택토 2판을 해보세요''로 시작하여 자유 발언을 이끌어낸 뒤, 시뮬레이터를 프로젝터에 띄우고 [\emoji{play-button} 학습 시작]을 클릭합니다. 3,000\textasciitilde5,000판 학습 후 [\emoji{bust-in-silhouette} 사람 대전]으로 전환하여 교사가 일부러 지는 시연을 합니다.

\begin{tipbox}[\emoji{light-bulb} 포인트]
학생이 ``AI가 처음엔 바보 $\to$ 점점 똑똑해진다''는 변화를 직접 목격하게 합니다.
\end{tipbox}

\subsubsection{\emoji{books} 개념 (8분) — 강화학습 3요소}

\begin{center}
\renewcommand{\arraystretch}{1.4}
\begin{tabular}{|l|l|}
\hline
\rowcolor{tableheader}
\textbf{개념} & \textbf{비유} \\
\hline
에이전트 & 처음 틱택토를 배우는 학생 \\
\hline
Q-테이블 & 학생의 비밀 노트 (``여기 두면 좋았어!'') \\
\hline
보상 & 이기면 사탕, 지면 벌칙 \\
\hline
탐험($\varepsilon$) & ``오늘은 새로운 수를 한번 둬볼까?'' \\
\hline
활용($1-\varepsilon$) & ``저번에 이겼던 그 수를 또 쓰자!'' \\
\hline
\end{tabular}
\end{center}

\subsubsection{\emoji{magnifying-glass-tilted-left} 탐색 (10분) — 교사 주도 관찰 (활동지 1)}

학생이 [\emoji{play-button} 학습 시작]을 클릭하고, 500판 $\to$ 5,000판 $\to$ 20,000판 순서로 정지하며 단계 이름, 승률 그래프 모양, $\varepsilon$ 값, Q-상태 수를 활동지 1에 기록합니다.

\subsubsection{\emoji{test-tube} 실습 (14분) — 학생 주도 실험 (활동지 2, 3)}

\textbf{활동 A: 파라미터 실험 (활동지 2)} — 초기화 $\to$ $\alpha$를 0.01로 $\to$ 10,000판 학습 $\to$ 다시 초기화 $\to$ $\alpha$를 0.8로 $\to$ 10,000판 학습 $\to$ 비교 기록.

\textbf{활동 B: AI와 대전 (활동지 3)} — 20,000판 이상 학습된 AI에서 사람 대전, Q-값 히트맵 관찰, AI 판단 근거 기록.

\subsubsection{\emoji{memo} 정리 (8분) — 성찰 및 공유}

학생 2\textasciitilde3명이 실험 결과를 발표하고, 교사가 핵심 개념(에이전트·환경·보상, Q-값, $\varepsilon$)을 정리합니다. ``AI가 스스로 배운다는 것은 어떤 의미일까요?'' 토론으로 마무리합니다.

% ============================================================
\section{학습 활동지}
% ============================================================

\subsection{활동지 1: AI 성장 관찰 일지}

\textbf{목표}: AI가 학습하면서 성장하는 과정을 단계별로 관찰하고 기록한다.

\begin{center}
\renewcommand{\arraystretch}{1.6}
\begin{tabular}{|c|c|c|p{3.5cm}|c|c|}
\hline
\rowcolor{tableheader}
\textbf{관찰 시점} & \textbf{단계 이름} & \textbf{등급} & \textbf{승률 그래프 모양} & \textbf{$\varepsilon$ 값} & \textbf{Q-상태 수} \\
\hline
500판 & (\hspace{1.5cm}) & (\hspace{0.8cm}) & & (\hspace{0.8cm})\% & (\hspace{0.8cm})개 \\
\hline
5,000판 & (\hspace{1.5cm}) & (\hspace{0.8cm}) & & (\hspace{0.8cm})\% & (\hspace{0.8cm})개 \\
\hline
20,000판 & (\hspace{1.5cm}) & (\hspace{0.8cm}) & & (\hspace{0.8cm})\% & (\hspace{0.8cm})개 \\
\hline
\end{tabular}
\end{center}

\textbf{생각해보기:}

\begin{enumerate}
\item 에피소드가 늘어나면서 가장 눈에 띄게 변한 것은 무엇인가요?
\item 무승부 비율은 어떻게 변했나요? 그 이유를 추측해보세요.
\item $\varepsilon$ 값이 줄어들면 AI의 행동은 어떻게 달라질까요?
\end{enumerate}

\subsection{활동지 2: 파라미터 실험 — 학습률($\alpha$)의 효과}

\textbf{목표}: 학습률($\alpha$)이 AI 학습에 미치는 영향을 실험으로 확인한다.

\begin{center}
\renewcommand{\arraystretch}{1.6}
\begin{tabular}{|l|c|c|c|}
\hline
\rowcolor{tableheader}
\textbf{항목} & \textbf{실험 A ($\alpha = 0.01$)} & \textbf{실험 B ($\alpha = 0.30$)} & \textbf{실험 C ($\alpha = 0.80$)} \\
\hline
10,000판 후 등급 & (\hspace{1.5cm}) & (\hspace{1.5cm}) & (\hspace{1.5cm}) \\
\hline
강도 (0\textasciitilde100) & (\hspace{1.5cm}) & (\hspace{1.5cm}) & (\hspace{1.5cm}) \\
\hline
그래프 안정적? & 예 / 아니오 & 예 / 아니오 & 예 / 아니오 \\
\hline
Q-상태 수 & (\hspace{1.5cm})개 & (\hspace{1.5cm})개 & (\hspace{1.5cm})개 \\
\hline
\end{tabular}
\end{center}

\begin{tipbox}[\emoji{light-bulb} 비유]
학습률은 ``한 번에 얼마나 많이 고칠까?''입니다. 시험을 보고 오답노트를 할 때, 너무 조금 고치면 실력이 안 늘고, 너무 많이 고치면 맞았던 것도 헷갈려요!
\end{tipbox}

\subsection{활동지 3: AI와 대전 — Q-값 히트맵 읽기}

\textbf{목표}: AI가 어떤 기준으로 수를 선택하는지 Q-값 히트맵을 해석한다.

AI와 3판을 진행하면서, 매 판마다 내가 둔 첫수, AI가 둔 첫수, Q-값 히트맵에서 가장 진한 초록색 칸, 빨간색 칸 유무, 대국 결과를 기록합니다.

\begin{center}
\renewcommand{\arraystretch}{1.4}
\begin{tabular}{|l|l|}
\hline
\rowcolor{tableheader}
\textbf{색상} & \textbf{의미} \\
\hline
진한 초록 & AI가 ``여기에 두면 이길 확률이 높아!'' \\
\hline
연한 초록 & ``나쁘지 않은 수'' \\
\hline
빨간색 & ``여기 두면 불리해...'' \\
\hline
회색 & 이미 돌이 놓인 자리 (선택 불가) \\
\hline
\end{tabular}
\end{center}

% ============================================================
\section{학습 질문}
% ============================================================

\subsection{초급 (이해) — ``무엇인가요?''}

\begin{center}
\renewcommand{\arraystretch}{1.4}
\begin{tabular}{|c|p{7cm}|p{4cm}|}
\hline
\rowcolor{tableheader}
\textbf{\#} & \textbf{질문} & \textbf{기대 답변 키워드} \\
\hline
1 & 강화학습의 3가지 핵심 요소는 무엇인가요? & 에이전트, 환경, 보상 \\
\hline
2 & Q-테이블은 무엇을 저장하는 표인가요? & 상태별 행동의 가치 (Q-값) \\
\hline
3 & 이 시뮬레이터에서 이기면 몇 점을 받나요? & $+1.0$ \\
\hline
4 & 에피소드(episode)란 무엇인가요? & 한 판의 게임 \\
\hline
5 & Q-값 히트맵에서 진한 초록색은 무슨 뜻? & AI가 좋다고 판단하는 위치 \\
\hline
\end{tabular}
\end{center}

\subsection{중급 (적용) — ``어떻게 달라지나요?''}

\begin{center}
\renewcommand{\arraystretch}{1.4}
\begin{tabular}{|c|p{6cm}|p{5cm}|}
\hline
\rowcolor{tableheader}
\textbf{\#} & \textbf{질문} & \textbf{실험 방법} \\
\hline
1 & $\alpha$를 0.01로 낮추면 AI 성장 속도가? & $\alpha = 0.01$ vs $0.30$ 비교 \\
\hline
2 & $\gamma$를 0.1로 낮추면 전략이? & $\gamma = 0.1$ vs $0.9$ 비교 \\
\hline
3 & $\varepsilon$ 감소율을 0.99로 빠르게 줄이면? & 감소율 $0.99$ vs $0.9995$ 비교 \\
\hline
4 & ``X 단독''과 ``자기대국''의 차이? & 학습 방식 전환 후 비교 \\
\hline
5 & 500판 AI vs 50,000판 AI 차이가 느껴지나요? & 사람 대전에서 체감 \\
\hline
\end{tabular}
\end{center}

\subsection{고급 (분석) — ``왜 그런가요?''}

\begin{enumerate}
\item AI가 충분히 학습하면 왜 무승부가 가장 많아지나요?
\item 대칭성 정규화를 하면 왜 학습 효율이 7\textasciitilde8배 올라가나요?
\item 탐험($\varepsilon$)을 처음부터 0\%로 설정하면 왜 AI가 잘 배우지 못할까요?
\item 역방향 갱신이 왜 순방향보다 효과적인가요?
\item 무승부를 0이 아닌 $+0.3$으로 설정한 이유는?
\end{enumerate}

\subsection{확장 (창의) — ``어디에 쓸 수 있을까요?''}

\begin{enumerate}
\item 강화학습으로 배울 수 있는 다른 게임에는 무엇이 있을까요?
\item 로봇이 걷는 법을 배울 때, 에이전트·환경·보상은 각각 무엇일까요?
\item 유튜브 추천 알고리즘도 강화학습과 비슷한 점이 있을까요?
\item 초등학생에게 강화학습을 설명한다면 어떤 비유를 사용할 건가요?
\item AI가 스스로 배운다면, AI에게 ``가르침''이 필요 없을까요?
\end{enumerate}

% ============================================================
\section{수업지도안}
% ============================================================

\subsection{수업 정보}

\begin{center}
\renewcommand{\arraystretch}{1.4}
\begin{tabular}{|l|l|}
\hline
\rowcolor{tableheader}
\textbf{항목} & \textbf{내용} \\
\hline
과목 & 컴퓨팅 사고와 AI 교육 \\
\hline
대상 & 초등 예비교사 (대학 2\textasciitilde4학년) \\
\hline
차시 & 1차시 (45분) \\
\hline
단원 & 머신러닝의 이해 — 강화학습 \\
\hline
주제 & AI가 게임을 배우는 방법: Q-Learning과 강화학습 \\
\hline
\end{tabular}
\end{center}

\subsection{성취 기준}

\begin{center}
\renewcommand{\arraystretch}{1.4}
\begin{tabular}{|l|p{9cm}|}
\hline
\rowcolor{tableheader}
\textbf{영역} & \textbf{성취 기준} \\
\hline
지식 & 강화학습의 3요소와 Q-Learning의 기본 원리를 설명할 수 있다 \\
\hline
기능 & 파라미터를 조절하여 AI 학습 결과를 비교·분석할 수 있다 \\
\hline
태도 & AI의 학습 과정을 관찰하며 탐구적 태도를 기르고, AI 교육의 가치를 인식한다 \\
\hline
\end{tabular}
\end{center}

\subsection{평가 계획}

\begin{center}
\renewcommand{\arraystretch}{1.4}
\begin{tabular}{|l|p{5cm}|l|}
\hline
\rowcolor{tableheader}
\textbf{평가 영역} & \textbf{평가 내용} & \textbf{평가 방법} \\
\hline
지식 이해 & 강화학습 3요소를 올바르게 설명 & 활동지 1 서술 \\
\hline
탐구 능력 & 파라미터 변경 후 결과 비교·분석 & 활동지 2 기록 \\
\hline
관찰·해석 & Q-값 히트맵을 올바르게 해석 & 활동지 3 기록 \\
\hline
의사소통 & 실험 결과를 논리적으로 발표 & 발표 관찰 \\
\hline
\end{tabular}
\end{center}

% ============================================================
\section{교사용 팁}
% ============================================================

\subsection{수업 운영 팁}

\begin{center}
\renewcommand{\arraystretch}{1.4}
\small
\begin{tabular}{|p{3.5cm}|p{8cm}|}
\hline
\rowcolor{tableheader}
\textbf{상황} & \textbf{대처 방법} \\
\hline
학생 PC가 느려서 학습이 오래 걸림 & 속도 슬라이더를 1ms로 낮추면 배치 크기 500으로 빠르게 진행 \\
\hline
그래프가 심하게 요동침 & 정상 현상임을 설명. 50판 평균이므로 초기에는 변동이 큼 \\
\hline
``왜 AI가 저한테 져요?'' & 학습량이 부족한 경우. 20,000판 이상 학습 필요 안내 \\
\hline
초기화를 깜빡하고 실험함 & [\emoji{clipboard} 보고서] 다운로드로 현재 상태 기록 후 초기화 \\
\hline
파라미터를 원래대로 못 돌림 & 기본값 안내: $\alpha=0.30$, $\gamma=0.90$, $\varepsilon$ 감소율$=0.9995$ \\
\hline
\end{tabular}
\end{center}

\subsection{자주 나오는 학생 질문}

\begin{center}
\renewcommand{\arraystretch}{1.4}
\small
\begin{tabular}{|p{4cm}|p{7.5cm}|}
\hline
\rowcolor{tableheader}
\textbf{질문} & \textbf{답변 가이드} \\
\hline
``AI끼리 대전하면 누가 이겨요?'' & 충분히 학습하면 거의 무승부. [\emoji{crossed-swords} AI 대전]으로 직접 확인 \\
\hline
``다른 게임에도 적용할 수 있나요?'' & 상태가 유한한 게임에 적합. 바둑은 딥러닝 필요 \\
\hline
``왜 무승부에 +0.3을 주나요?'' & 0이면 방어적 플레이를 못 배움. 약간의 보상으로 유도 \\
\hline
\end{tabular}
\end{center}

% ============================================================
\section{수업 확장 아이디어}
% ============================================================

\subsection{심화 활동}

\begin{center}
\renewcommand{\arraystretch}{1.4}
\begin{tabular}{|l|p{6cm}|c|}
\hline
\rowcolor{tableheader}
\textbf{활동} & \textbf{내용} & \textbf{소요 시간} \\
\hline
AI 대전 토너먼트 & 다른 파라미터로 학습한 AI끼리 대전 & 15분 \\
\hline
학습 보고서 비교 & 보고서 다운로드 후 조별 비교 발표 & 10분 \\
\hline
보상 체계 토론 & ``무승부를 $-0.5$로 바꾸면?'' 예측·토론 & 10분 \\
\hline
X vs O 비교 & X 단독 vs O 단독 학습 후 성능 비교 & 15분 \\
\hline
\end{tabular}
\end{center}

\subsection{교과 연계}

\begin{center}
\renewcommand{\arraystretch}{1.4}
\begin{tabular}{|l|p{9cm}|}
\hline
\rowcolor{tableheader}
\textbf{교과} & \textbf{연계 내용} \\
\hline
수학 & 확률과 통계 — 승률 그래프 읽기, 평균 개념 \\
\hline
도덕 & AI 윤리 — ``AI가 스스로 배운다면 행동 책임은 누구에게?'' \\
\hline
국어 & 설명문 쓰기 — ``AI가 틱택토를 배우는 과정'' 설명문 \\
\hline
실과 & 소프트웨어 교육 — 알고리즘 개념, 컴퓨팅 사고력 \\
\hline
\end{tabular}
\end{center}

\subsection{다른 시뮬레이터와의 비교 수업}

\begin{center}
\renewcommand{\arraystretch}{1.4}
\begin{tabular}{|l|p{8cm}|}
\hline
\rowcolor{tableheader}
\textbf{비교 대상} & \textbf{비교 포인트} \\
\hline
사과 가격 예측 (선형 회귀) & 지도학습 vs 강화학습 — 정답 유무의 차이 \\
\hline
KNN 분류 & 데이터 기반 vs 경험 기반 — 학습 방식의 차이 \\
\hline
K-means 군집화 & 비지도학습 vs 강화학습 — 목표의 차이 \\
\hline
\end{tabular}
\end{center}

% ============================================================
\newpage
\section*{참고자료}
% ============================================================

\begin{center}
\renewcommand{\arraystretch}{1.4}
\begin{tabular}{|c|l|l|l|}
\hline
\rowcolor{tableheader}
\textbf{\#} & \textbf{자료명} & \textbf{유형} & \textbf{비고} \\
\hline
1 & rl-tictactoe-v3.html & 시뮬레이터 & 소스 코드 \\
\hline
2 & RL틱택토\_사용설명서\_v1.0.0.0 & 프로젝트 문서 & 전체 기능 상세 \\
\hline
3 & RL틱택토\_이론해설서\_v1.0.0.0 & 프로젝트 문서 & Q-Learning 원리 \\
\hline
4 & RL틱택토\_퀵가이드\_v1.0.0.0 & 프로젝트 문서 & 빠른 시작 참조 \\
\hline
5 & 코딩관리지침\_v3.0.0.0 & 프로젝트 지침 & §7 수업 활용 자료 \\
\hline
6 & 문서작성지침서\_v1.6.0.0 & 프로젝트 지침 & 톤/어조 \\
\hline
\end{tabular}
\end{center}

\section*{생성/수정 이력}

\begin{center}
\renewcommand{\arraystretch}{1.3}
\small
\begin{tabular}{|l|l|l|l|p{5cm}|l|}
\hline
\rowcolor{tableheader}
\textbf{버전} & \textbf{날짜} & \textbf{시간} & \textbf{변경 수준} & \textbf{변경 내용} & \textbf{작성자} \\
\hline
v1.0.0.0 & 2026-02-14 & 22:30 & 최초 & 초안 작성 — 수업 시나리오, 활동지 3종, 학습 질문 4단계, 수업지도안, 교사용 팁, 확장 아이디어 & Claude \\
\hline
\end{tabular}
\end{center}

\vfill
\begin{center}
\small\color{lightbrown}
\emph{RL 틱택토 활용가이드 — 교육용 Q-Learning 시뮬레이터의 수업 적용 가이드}
\end{center}

\end{document}
